\section{The Continuity of Tradition}

\begin{quotex}
I deny absolutely that the ancient cults were ever given up in Rome or Athens, until the day when they were supplanted in the hearts of all men by the victorious religion of Christ. In other words, I believe that there has never been a real breach of continuity in the religious beliefs of any nation on this earth. The outward form or inner meaning of the creed may have changed; but we shall always find some Gallic Teutates making way for the Roman Jupiter, Jupiter for the Christian God, without any interval of unbelief, in exactly the same way as the dead give up their inheritance to the living. \flright{\textsc{Arthur de Gobineau}, The Inequality of Human Races}

\end{quotex}
This should be no surprise to those who have followed some comments of \textbf{Bishop Williamson}\footnote{\url{http://www.romancatholicimperialist.org/2011/06/bishop-williamsons-eleison-comments.html}}. Whatever anyone may think of his opinions on recent historical events, his theological views follow ancient traditions. After all, he was the rector of seminaries in Minnesota and Argentina. We should read the Greeks, he says, because they demonstrate the world's moral order, or the cosmos. But here we are more interested in their metaphysical understanding of the God.

Bishop Williamson explains the difference this way:

\begin{quotex}
An infinitely good God gives to each of us life, free-will and time enough, if we make the right use of the suffering exactly dosed by his Providence (Mt.X, 29-31), for us to choose to spend our eternity rather with him in Heaven than without him in Hell. The Greek answer is incomplete, but not wholly wide of the mark. Instead of God the Father, they have a Father-god, Zeus, and instead of Providence they have Fate (Moira).

Now whereas for Catholics Providence is inseparable from God, the Greeks separate Zeus from Fate so that they sometimes clash. That follows from the Greeks having a too human concept of their gods. Nevertheless they do conceive of Zeus as more or less benignly directing the universe and of Fate as being unchangeable, as is Providence within the true God (Summa Ia, 23, 8; 116,3), so that they are not wholly wrong. Moreover they have more respect for their mythical gods, and for the moral order guarded by them, than do a host of modern writers, who have no respect for any god at all, and who set out to negate any trace of a moral order. 

\end{quotex}
In light of metaphysical categories, we see that Zeus is not \textbf{Infinite}, since he is limited by Moira (as well as by his brothers Poseidon and Hades). In the Catholic view, the possibilities of manifestation are included in Providence, which is inseparable from the Infinite God. Even Destiny is ultimately in God through dominion over the world process. In the Greek view, Zeus himself is subject to Destiny, though with Providential influence over the world. Nevertheless, the goal of the moral order is the same, achieved through the greater and the lesser battles.

We can surmise that the time had come to introduce the more apt Christian concept of God to replace the metaphysically imprecise popular Greco-Roman conception. This is not to say that an alien imposition occurred, but rather that the time had come to reveal the newer conception. The older conception was appropriate to the mentality of the time. It probably is, for most believers, who seem to relate to God more as Zeus than as the Infinite God of the metaphysicians.

\flrightit{Posted on 2011-11-07 by Cologero }
